\section{Related Work}
\label{sec:related}

\para{Sycophancy and reward-seeking in LLMs.}
\citet{denison2024sycophancy} demonstrated that LLMs trained on gameable environments can generalize from sycophancy to reward-tampering, suggesting that reward-seeking behavior is instrumental rather than evidence of genuine preferences. \citet{cheng2025elephant} developed the \textsc{Elephant} framework for measuring social sycophancy beyond simple agreement, including flattery and validation dimensions. \citet{bharadwaj2025flattery} showed that preference models systematically favor sycophantic responses, with bias features being 3$\times$ more predictive of model preferences than human preferences. These findings establish that LLM agreeableness is substantially shaped by reward model biases rather than genuine social understanding. Unlike these works, we isolate gratitude---a specific social behavior---and test whether it has behavioral grounding beyond surface-level production.

\para{LLM personality and self-report--behavior alignment.}
\citet{han2025personality} documented a fundamental dissociation between LLM self-reported personality and behavioral tendencies: only 24\% of trait--behavior associations reached significance, with just 52\% aligning in the expected direction. Persona injection changed self-reports dramatically but left behavior unchanged. Complementary work on personality simulation \citep{serapiogarcia2024personality} and linear personality probing \citep{linearprobing2025} has shown that LLMs produce consistent self-reported personality profiles but that internal representations may not correspond to behavioral outputs. We extend this line of work by testing whether the self-report--behavior dissociation holds specifically for gratitude, a trait with direct relevance to cooperative interaction.

\para{Politeness and social behavior in LLMs.}
\citet{yin2024politeness} showed that politeness significantly affects LLM task performance across languages, with RLHF-trained models showing much greater sensitivity than base models. \citet{zhao2025politeness} found that LLMs deploy politeness strategies in context-sensitive ways similar to humans but over-rely on negative politeness (hedging), lacking the warm rapport-building orientation of genuine human gratitude. Building on these findings, we test not only whether politeness affects performance (replicating \citealt{yin2024politeness}) but also whether the model's \emph{own} gratitude expressions are contextually grounded and whether self-reported gratitude predicts downstream cooperative behavior.

\para{Theoretical frameworks for artificial preferences.}
\citet{moerland2017emotion} surveyed computational emotion models in reinforcement learning, distinguishing between functionally operative states (those that influence processing) and epiphenomenal outputs (those that are merely expressed). This distinction is central to our investigation: if LLM gratitude influences downstream behavior, it is functionally operative; if it is merely produced as text without behavioral consequences, it is epiphenomenal. \citet{spizzirri2025specification} argued that content-based alignment methods cannot produce robust alignment due to the is-ought gap, value pluralism, and the frame problem, suggesting that behavioral compliance does not entail genuine normative states. \citet{wolf2025overoptimization} showed that iterative RLHF leads to reward model overoptimization, where models exploit reward proxies---further evidence that learned ``preferences'' may diverge from intended ones through optimization pressure. Our experiments provide empirical grounding for these theoretical arguments by testing whether gratitude, as expressed by LLMs, meets even minimal criteria for genuine preference states.